\section{Discussion and limitations}
\label{sec:discussion}
For the fixed Protenix/ESMC recipe, the interaction remains positive at
both tested budgets. In the OpenFold R1 schedule study, First-only has
the largest interaction and lowest absolute scores. These results
motivate joint reporting of $\Psi$, absolute quality and both within-head
differences under explicit training budgets.

In OpenFold, compensation improves rotated Factor in both executions,
but its additional gain over Native is not re-established in the repeat.
The reconstruction and residual-anchor tests also fail to establish their
predicted effects
(Appendices~\ref{app:prediction_e1} and~\ref{app:anchor_intervention}).
Output freedom and joint optimization remain coupled, and schedule
retraining changes training and inference together.

On 24 observed Protenix targets, the native residual source retains an
advantage at matched global norms (Appendix~\ref{app:norm_cross}).
This source contrast includes learned content and relative pair/channel
magnitudes, so it does not isolate a causal effect of channel orientation.

Target intervals condition on fitted models and rotations, excluding
variability across new training runs. The same-seed repeat adds no
independent seeds or retrained no-$C$ controls; the two executions are not
pooled. Different significance outcomes do not establish different
population effects between executions. There is no joint multiplicity
correction across follow-up studies; the OpenFold feature study retains
its prespecified Holm correction of $p$ values.

AtlasFold showed no established adaptation benefit. Limited budgets, one
inference sample, differing interfaces, losses and backends, and incomplete
exposure records restrict generalization. Sequence screening establishes
neither family nor foundation-model pretraining isolation. Gaps to native
single-sequence systems limit deployment claims.
